%% ch01-introduction.tex — JA4proxy Operator's User Guide
%% Chapter 1: Introduction

\chapter{Introduction}
\index{introduction}

\newthought{Most TLS security products work by breaking encryption.}
They sit between the client and the server, decrypt the traffic, inspect it,
re-encrypt it, and forward it on. This gives them visibility into HTTP content,
but it comes at significant cost: the product must hold your TLS private keys,
your users are exposed to an additional certificate chain they did not negotiate,
and performance overhead is substantial.

\ja\ takes a different approach. It never decrypts anything.

\section{The Core Insight: TLS Fingerprinting}
\index{TLS fingerprinting}
\index{JA4 fingerprint}

The TLS handshake begins with a \emph{ClientHello} message sent in plaintext before
any encryption is established. This message contains a structured description of the
client's TLS capabilities: which protocol versions it supports, which cipher suites
it offers, which extensions it uses, and in what order.

The combination of these fields is highly characteristic of the software that generated
them. Chrome 120 on Windows produces a different ClientHello than curl 8.3, which
produces a different ClientHello than a Cobalt Strike beacon, which produces a different
ClientHello than an automated credential-stuffing tool.

The \emph{JA4} algorithm\sidenote{JA4 is a fingerprinting algorithm by FoxIO. Specification
at \url{https://github.com/FoxIO-LLC/ja4}} hashes these ClientHello fields into a compact,
human-readable fingerprint string. A typical fingerprint looks like
\fingerprint{t13d1113h2\_8daaf6152771\_b0da82dd1658}. The prefix encodes the
TLS version, the middle section encodes the cipher and extension set, and the suffix
encodes the ALPN values.

\ja\ computes this fingerprint for every incoming connection — before the handshake
completes, before any HTTP request is made — and uses it as the primary signal for
traffic classification.

\begin{notebox}[The architectural guarantee]
Because \ja\ reads only the ClientHello plaintext, it never holds TLS private keys
and never terminates TLS. The encrypted session passes through unchanged. Your backend
server terminates TLS exactly as it would without \ja\ in the path.
\end{notebox}

\section{What JA4proxy Does Not Do}
\index{non-goals}

Understanding the boundaries is as important as understanding the capabilities.

\ja\ \emph{does not}:
\begin{itemize}
  \item Inspect HTTP request content, headers, or bodies
  \item Decrypt TLS traffic
  \item Parse application-layer protocols
  \item Replace a WAF for content-based attack detection (SQLi, XSS, and similar)
  \item Authenticate individual users
\end{itemize}

\ja\ \emph{does}:
\begin{itemize}
  \item Identify and block known-bad TLS fingerprints (C2 tools, crawlers, bots)
  \item Enforce minimum TLS version and cipher suite requirements
  \item Detect and rate-limit automated scanning behaviour
  \item Block connections from known-bad IP addresses, ASNs, and countries
  \item Identify Tor exit nodes, datacenter IP ranges, and VPN providers
  \item Detect beaconing patterns (regular-interval connections characteristic of C2 traffic)
  \item Apply reputation data from AbuseIPDB and Spamhaus DROP lists
\end{itemize}

\section{The Network Architecture}
\index{architecture}

\ja\ sits between your load balancer and your backend. In the reference deployment:

\begin{figure}[ht]
\begin{lstlisting}[language=bash, numbers=none, backgroundcolor=\color{white}, frame=none, xleftmargin=0pt, framexleftmargin=0pt]
Internet
  |
  v  :443 (TLS, PROXY protocol v2)
HAProxy (load balancer)
  |
  v  :8080 (TCP passthrough)
JA4proxy x N  <----> Redis (shared state)
  |
  v  :443 (TLS, unmodified)
Backend server
\end{lstlisting}
\end{figure}

HAProxy terminates nothing — it operates in TCP passthrough mode, forwarding the raw
TLS stream to \ja\ with a PROXY protocol v2 header prepended so the proxy knows the
real client IP address. \ja\ reads the ClientHello, makes its decision, and either
forwards the connection to the backend unchanged or closes it.

Multiple \ja\ instances share state through Redis: the JA4 blacklist and whitelist,
IP bans, rate-limiting windows, and enrichment results are all available to every
instance.\sidenote{This means you can scale to many proxy instances without
any sticky-session requirement.}

\section{The Decision Pipeline}
\index{pipeline}
\index{bypass!fast path}

Every connection passes through a two-stage pipeline.

\subsection{Stage 1: Fast-Path Bypasses}

Before any scoring occurs, a set of conditions can short-circuit the pipeline
entirely. These are \emph{bypass checks} — they produce an immediate ALLOW or BLOCK
decision without consulting the scorer.

\begin{table}[ht]
  \small
  \begin{tabularx}{\linewidth}{lXl}
    \toprule
    \textbf{Condition} & \textbf{Behaviour} & \textbf{Configurable} \\
    \midrule
    h2/h1 ALPN (browser) & \badgeallow{} immediately & Yes \\
    JA4 in whitelist & \badgeallow{} immediately & Yes \\
    Valid mTLS client cert & \badgeallow{} immediately & Yes \\
    IP in static allowlist & \badgeallow{} immediately & Yes \\
    \midrule
    JA4 in blacklist & \badgeblock{} immediately (RST) & Yes \\
    Country in blocklist & \badgeblock{} immediately & Yes \\
    Spamhaus DROP match & \badgeblock{} immediately & Yes \\
    TLS 1.0/1.1 version & \badgeblock{} immediately (RST) & Yes \\
    \bottomrule
  \end{tabularx}
  \caption{Fast-path bypass conditions and their defaults}
\end{table}

The browser ALPN bypass deserves special mention. Any connection advertising \code{h2}
or \code{h1} (standard browser ALPN values) is allowed immediately, regardless of
any other rule. This is an architectural guarantee against false positives — legitimate
browser traffic can never be blocked by \ja, regardless of configuration.\sidenote{This
guarantee exists because the cost of blocking a real user is always higher than the cost
of letting through a bot that happens to advertise browser ALPN values. The proxy is
designed around this asymmetry.}

\subsection{Stage 2: Signal Collection and Scoring}

Connections that reach Stage 2 are analysed by a set of signal collectors running
in parallel. Each collector examines one aspect of the connection and produces a
\emph{RiskSignal} — a numeric score contribution and a label. The signals are:

\begin{table}[ht]
  \small
  \begin{tabularx}{\linewidth}{lX}
    \toprule
    \textbf{Signal} & \textbf{What it examines} \\
    \midrule
    TLS enforcer & TLS version, cipher strength \\
    SNI analyser & Missing SNI, IP-literal, DGA-pattern hostnames \\
    TCP analyser & JA4T fingerprint, connection lifespan, concurrency \\
    ASN classifier & Datacenter, Tor exit, VPN ranges \\
    FCrDNS enrichment & Forward-confirmed reverse DNS; residential classification \\
    Blocklist check & Spamhaus DROP/EDROP match (as signal, not bypass) \\
    Beaconing detector & IAT coefficient of variation; regular-interval connections \\
    AbuseIPDB & IP reputation score from community abuse reports \\
    RDAP enrichment & Org-level reputation; netblock classification \\
    Analytics findings & Cross-instance campaign and slow-scan detection \\
    \bottomrule
  \end{tabularx}
  \caption{Signal collectors in Stage 2}
\end{table}

The composite scorer aggregates all RiskSignals into a single score from 0 to 100.
The action decider then maps that score to a decision based on the current \emph{dial} setting.

\section{The Dial}
\index{dial}
\index{monitor mode}

The dial is a single number from 0 to 100 that controls how aggressively \ja\ acts
on risk scores.

\begin{description}[leftmargin=2.5cm]
  \item[\dialzero] The proxy scores all traffic but takes no blocking action.
    Every connection is allowed. Events are logged and metrics are recorded.
    This is the default and the correct setting for initial deployment.
  \item[\dialhundred] The proxy applies its configured thresholds. Connections
    with scores above the block threshold are blocked; those above the tarpit
    threshold are tarpitted; those above the rate-limit threshold are rate-limited.
\end{description}

\begin{importantbox}[Why the default is zero]
The default \code{dial=0} is deliberate. On first deployment, you do not yet know
what your legitimate traffic looks like. Deploying with blocking enabled risks
immediately breaking real users. Spend time at dial=0 reviewing the logs and Grafana
dashboard. Only raise the dial when you are confident the scoring matches reality.
\end{importantbox}

Values between 0 and 100 proportionally raise the thresholds at which each action
triggers. The dial is designed to be raised gradually: from 0 to 25 (flag only),
then to 50 (rate limiting), then to 75 (tarpit high-risk), then to 100 (full blocking).
\emph{See Chapter~\ref{ch:configuration} for the full dial rollout procedure.}

\section{Self-Healing Design}
\index{self-healing}
\index{ban!TTL}

Every ban in \ja\ has a 5-minute TTL.\sidenote{This is intentional and cannot be
configured away. The rationale: a 5-minute ban on a scanner costs them the same
operational disruption as a permanent ban, because scanners retry anyway. But a
5-minute window limits the blast radius when a legitimate IP is incorrectly flagged.}
There are no permanent bans by default. The system is designed to fail open: if
Redis is unreachable, the proxy allows connections rather than blocking them.

This approach accepts a small number of false negatives (bad traffic slipping through
during a recovery window) in exchange for a strong guarantee against false positives
(real users being locked out).

\section{What This Guide Covers}

This guide is the operational reference: how to install, configure, operate, monitor,
and troubleshoot \ja. It assumes you have a working knowledge of Docker, TLS basics,
and network security concepts.

For architecture decision records, signal algorithm specifications, full Prometheus
metric definitions, Redis schema documentation, and API references, see the
\emph{Technical Reference Manual}.
